% This is samplepaper.tex, a sample chapter demonstrating the
% LLNCS macro package for Springer Computer Science proceedings;
% Version 2.21 of 2022/01/12
%
\documentclass[runningheads]{llncs}

\usepackage[T1]{fontenc}
% T1 fonts will be used to generate the final print and online PDFs,
% so please use T1 fonts in your manuscript whenever possible.
% Other font encondings may result in incorrect characters.
%

\usepackage{graphicx}
\usepackage{epsfig}
\usepackage{subcaption}
\usepackage{calc}
\usepackage{amssymb}
\usepackage{amstext}
\usepackage{amsmath}
\usepackage{multicol}
\usepackage{pslatex}
\usepackage{algorithm2e}
\usepackage[bottom]{footmisc}
\usepackage{booktabs}
\usepackage{multirow}
\usepackage{svg}

% Used for displaying a sample figure. If possible, figure files should
% be included in EPS format.
%
% If you use the hyperref package, please uncomment the following two lines
% to display URLs in blue roman font according to Springer's eBook style:
%\usepackage{color}
%\renewcommand\UrlFont{\color{blue}\rmfamily}
%\urlstyle{rm}
%
\begin{document}
%
\title{PAMR: Link Prediction on Knowledge Graphs via Path-Aware Multi-modal Reasoning}
%
\titlerunning{Link Prediction on KG via Path-Aware Multi-modal Reasoning}
% If the paper title is too long for the running head, you can set
% an abbreviated paper title here
%
\author{Hai-Dang Huynh-Lam\inst{1,2}\orcidID{0000-0001-9487-7281} \and
  Ngoc-Thao Nguyen\inst{1,2}\orcidID{1111-2222-3333-4444} \and
Ngoc-Thanh Le\inst{1,2}\orcidID{0009-0009-6426-2652}}
%
\authorrunning{Hai-Dang et al.}
% First names are abbreviated in the running head.
% If there are more than two authors, 'et al.' is used.
%
\institute{
  Faculty of Information Technology, University of Science, Ho Chi Minh City, Vietnam
  % \email{lncs@springer.com}\\
  % \url{http://www.springer.com/gp/computer-science/lncs}
  \and
  Vietnam National University, Ho Chi Minh City, Vietnam\\
\email{\{hlhdang,nnthao,lnthanh\}@fit.hcmus.edu.vn}}
%
\maketitle  % typeset the header of the contribution
%
\begin{abstract}
  The abstract should briefly summarize the contents of the paper in
  150--250 words.

  \keywords{Knowledge Graph \and Multi-modal Learning \and Path-based Reasoning \and Link Prediction \and Transformers.}
\end{abstract}
%
%
%

\section{Introduction}
\label{sec:introduction}

Knowledge Graphs (KGs) serve as foundational structures for organizing factual information into interconnected entities and relations. Link prediction is a fundamental task on KGs aimed at inferring missing connections to overcome inherent incompleteness. Traditional approaches predominantly focus on structural graph information to learn continuous vector representations for entities and relations. As real world data spans diverse formats, Multimodal Knowledge Graphs, or MMKGs, have emerged to incorporate auxiliary modalities such as images and text into the structural topology. While multimodal knowledge graph embedding methods successfully align these modalities to improve link prediction, they generally compress diverse information into a single unified entity representation. This fusion process obscures the explicit connectivity patterns and limits the ability to perform interpretable reasoning across heterogeneous modalities.

Path aware reasoning addresses this interpretability limitation by extracting and evaluating explicit structural sequences between source and target entities. However, existing path based models are primarily designed for concept only configurations and fail to exploit the rich intermediate signals provided by multimodal nodes. Furthermore, most architectures entangle the sequential modeling process with embedding initialization, which creates significant computational burdens when scaling to complex multimodal representations.

To address these challenges, we propose the Path-Aware Multimodal Reasoning framework, denoted as PAMR. This framework addresses link prediction on multimodal knowledge graphs by scoring structural paths between query entities using a Transformer encoder over frozen knowledge graph embedding token representations. By operating under a modality as entity setting, PAMR treats multimodal information as distinct graph nodes rather than static attachments. This allows the model to extract paths that route through shared visual or textual attributes, providing explicit and highly interpretable evidence for link prediction without the initialization deadlock common in continuous sequence modeling.

The core contributions of this work are summarized as follows.

\begin{itemize}
  \item First, to the best of our knowledge, this is the pioneer work proposing path aware reasoning for link prediction on MMKGs under a modality as entity setting, establishing a robust baseline utilizing the unified embeddings provided by VL-KGE.
  \item Second, the proposed path extraction process leverages both conceptual entities and multimodal entities, with experimental results demonstrating substantial performance gains derived from integrating these multimodal signals.
  \item Third, we introduce a Transformer based sequential architecture for scoring the extracted paths. This approach achieves results comparable to strong embedding baselines while offering significantly enhanced interpretability by tracing concrete multimodal reasoning trajectories.
\end{itemize}

\section{Related Work}
\label{sec:related-work}

The development of the PAMR framework lies at the intersection of several core research areas in artificial intelligence, including graph-based relational reasoning, multimodal representation learning, and sequential modeling. In this section, we survey the literature across three foundational dimensions to establish the methodological context and justify the necessity of PAMR.

\subsection{Path-Based Reasoning on Knowledge Graphs}
\label{sec:rw-path-reasoning}

Path-based reasoning on knowledge graphs extends conventional single-hop link prediction by extracting and evaluating multi-hop relational sequences that connect source and target entities. Foundational methods in this domain relied heavily on statistical random walks, such as the Path Ranking Algorithm (PRA)~\cite{lao2011random}. The transition to deep learning introduced architectures like NeuralLP~\cite{yang2017differentiable} and RNNLogic~\cite{qu2021rnnlogic}, which formulate path extraction as differentiable logic programming or probabilistic modeling, respectively, enabling end-to-end learning of logical rules. Concurrently, deep reinforcement learning (RL) was adapted to navigate graph spaces. DeepPath~\cite{xiong2017deeppath} and MINERVA~\cite{das2018go} model pathfinding as a Markov Decision Process, where an agent learns to step toward the target entity using shaped rewards. To mitigate the sparse reward issues common in long-path traversal, DIVA~\cite{chen2018variational} framed the problem as variational inference, while CURL~\cite{zhang2022learning} proposed a multi-agent hierarchical RL framework.

More recently, graph neural networks (GNNs) and sequence-to-sequence models have been employed to efficiently encode path structures. RED-GNN~\cite{zhang2022knowledge} and NBFNet~\cite{zhu2021neural} parameterize the Bellman-Ford algorithm with neural components to capture generalized path representations through message passing. AdaProp~\cite{zhang2023adaprop} improves upon this by adaptively filtering propagation paths, while SQUIRE~\cite{bai2022squire} introduces a Transformer-based sequence-to-sequence framework that translates queries directly into paths. However, most path-based methods rely on learned rather than frozen embeddings, use reinforcement learning or message-passing rather than Transformers, and do not handle multimodal entities. This interwining between embedding and the sequential modeling process increases computational burden and limiting the ability to leverage advanced representations.

\subsection{Multimodal Knowledge Graph Embeddings}
\label{sec:rw-mmkge}

Multimodal knowledge graph embedding (MMKGE) seeks to overcome the structural sparsity of traditional KGs by injecting auxiliary modalities, such as images and text, into the representation learning process. Early approaches, such as IKRL~\cite{xie2017image} and TransAE~\cite{wang2019transae}, demonstrated that aligning visual and textual features with structural embeddings could resolve entity ambiguity and improve link prediction. The MMKG dataset collection~\cite{liu2019mmkg} further established baselines for numeric and visual entity matching. Subsequent models focused on filtering modality noise and handling structural complexity; for instance, RSME~\cite{wang2021rsme} uses a forget-gate mechanism to discard conflicting multimodal signals, while OTE~\cite{tang2020orthogonal} maps multimodal entities via orthogonal relation transforms to preserve topological distances. Additionally, MMRNS~\cite{xu2022multimodal} incorporates cross-modal attention to enhance negative sampling with multimodal semantics.

State-of-the-art MMKGE architectures increasingly rely on Transformers. MKGformer~\cite{chen2022mkgformer} employs a hybrid multi-level fusion approach to integrate visual and textual representations. MoSE~\cite{zhao2022mose} introduces a modality-split framework to prevent interference during fusion, and VISTA~\cite{lee2023vista} extends multimodal encoding directly to the triplet level using specialized decoding Transformers. Finally, VL-KGE~\cite{efthymiou2026vlkge} leverages vision-language models to handle modality asymmetry in large-scale graphs. Notably, existing multimodal KGE methods fuse modalities into a single entity embedding but do not expose multimodal token sequences to a path-aware Transformer for reasoning, as the fusion primarily serves triple scoring rather than path discrimination. Therefore, no widely used benchmark treats multimodal information as entities on the graph that can connect to different conceptual entities rather than being bound to a single entity. This multimodal-as-entity setting is what the proposed method evaluates on.

\subsection{Transformers for Knowledge Graph Reasoning}
\label{sec:rw-transformers}

The application of Transformers and pre-trained language models to KGs leverages self-attention mechanisms to contextualize relational data. KG-BERT~\cite{yao2019kg} pioneered this by treating triples as natural language sequences for sequence classification. To address the computational bottleneck of cross-encoding every candidate, StAR~\cite{wang2021star} partitions triples into asymmetric context representations. Moving beyond purely textual encoding, CoKE~\cite{wang2019coke} tokenizes graph structures, passing edge and path sequences through a Transformer encoder to learn context-adaptive embeddings. HittER~\cite{chen2021hitter} extracts features through a hierarchical architecture, aggregating local neighborhood information using masked entity prediction.

For complex logical queries, KG-Transformer~\cite{liu2022mask} utilizes a masked pre-training strategy coupled with a Mixture-of-Experts sparse activation. Generative and prompt-based paradigms have also emerged, with KGT5~\cite{saxena2022kgt5} formulating link prediction as a sequence-to-sequence task, GenKGC~\cite{zhu2022genkgc} utilizing relation-guided prompts for target generation, and CSPrompter~\cite{chen2023csprompter} bridging the gap between symbolic logic and large language models via consistent structure prompting. However, Transformer-based KG methods typically encode triples or subgraphs rather than interleaved path sequences, use language-model pretraining rather than frozen KGE embeddings, and mostly target static text KGs rather than multimodal graphs, let alone the multimodal graphs under multimodal-as-entity setting.

% \subsection{Evaluation Approaches}
% \label{sec:rw-evaluation}

% Evaluation of link prediction models spans both training signal quality
% and test protocol rigor. Margin-based ranking losses, widely used in
% this domain, rely heavily on negative sampling quality. Traditional
% uniform random sampling~\cite{bordes2013translating} or Bernoulli
% distributions~\cite{wang2014knowledge} often yield easily
% distinguishable or false-negative samples, while self-adversarial
% weighting~\cite{sun2019rotate} and generative adversarial frameworks
% like KBGAN~\cite{cai2018kbgan} improve gradient informativeness at
% the cost of training stability. More efficient alternatives such as
% NSCaching~\cite{zhang2019nscaching} and CAKE~\cite{niu2022cake}
% achieve adversarial-level performance through cached hard negatives
% and commonsense priors, respectively. However, these methods are
% designed for triple-level discrimination and do not guarantee balanced
% positive-to-negative ratios per batch, making them unsuitable for
% path-level paired sampling.

% Standardized benchmarks complement these training-side advances by
% providing rigorous test protocols. WN18RR~\cite{dettmers2018convolutional}
% and FB15k-237~\cite{toutanova2015observed} eliminate test-set leakage
% and inverse relation bias, while
% YAGO3-10~\cite{mahdisoltani2013yago3} evaluates scalability across
% multilingual data. With the shift towards multimodality,
% domain-specific graphs like WikiArt-MKG~\cite{efthymiou2026vlkge}
% offer heterogeneous structures spanning artworks, artists, and textual
% attributes. Recent surveys~\cite{niu2025survey,liang2024survey} have
% cataloged these developments by task orientation and graph type.
% Nevertheless, standard benchmarks are either purely structural without
% multimodal features, or lack explicit path annotations for
% supervision, leaving WikiArt-MKG as the only dataset that combines
% both properties yet has not been used for Transformer-based path
% reasoning.

Taken together, existing work leaves several open questions that PAMR addresses: how to explicitly execute path-aware reasoning on heterogeneous multimodal entities, how to parse and project frozen complex-valued KGEs into a Transformer's token space without initialization deadlock, and how to stabilize sequence ranking via a perfectly balanced query-level paired sampling mechanism.

\section{Method}
\label{sec:method}
\subsection{Overview}
\label{sec:overview}

A KG is a directed multi-relational graph
$\mathcal{G} = (\mathcal{V}, \mathcal{E}, \mathcal{R})$ where
$\mathcal{V}$ is a set of entities, $\mathcal{R}$ a set of relation types,
and $\mathcal{E} \subseteq \mathcal{V} \times \mathcal{R} \times \mathcal{V}$
a set of labelled edges. Each edge $(s, r, o) \in \mathcal{E}$
denotes that subject entity $s$ is connected to object entity $o$
via relation $r$.

Path-based reasoning focuses on using structural paths and connection
patterns within KGs to infer new relations or answer queries. Given
a query $(s, r, ?)$, the target is to approximate the conditional
distribution $P(\hat{o} \mid s, r, \mathcal{G})$ over candidate
objects $\hat{\mathbb{O}}$. By using the shortest paths $p_o$ between $s$ and $\hat{o}$ in $\mathcal{G}$ as the most important factor contributing to the distribution,
the target reduces to $P(p_o \mid s, r)$. A scoring function
$\hat{f}_\theta$ is learned to predict this target and then assigns a real-valued score to each path given the
query relation $r$, afterward the candidate with the highest-scoring path is selected as the answer.

In a MMKG, entities are associated with heterogeneous modalities, for example artworks are inherently visual, while artists,
styles, genres, and locations are characterized textually. We model
this as a multimodal-as-attachment setting, where each entity $v \in \mathcal{V}$
possesses modality-specific embeddings $\mathbf{v}_i$ indexed by
$i \in \mathbb{M}$, with $\mathbb{M}$ denoting the modality set.
The unified embedding $\mathbf{v}$ is obtained
by aggregating over available modalities: $\mathbf{v} =
\text{AGGREGATE}_{i \in \mathbb{M}} (\mathbf{v}_i)$. In this proposal, we formulate the MMKG setting of modality-as-entity where each modality information is converted to a separate entity, this creates a mapping $\mathcal{M}: \mathcal{V} \to \mathbb{M}$ that maps each entity to its original modality.

The proposed method, PAMR, consists of four components as illustrated in Figure \ref{fig:overview}. First, a frozen
pre-trained KGE model supplies token embeddings for entities and relations.
Second, a bidirectional BFS extracts shortest paths between query pairs under the modality-as-entity setting.
Third, a Transformer encoder scores each path by contextualising
interleaved entity and relation tokens. Fourth, a margin-based ranking
loss trains the scorer on paired positive and negative paths. Each
component is detailed in the following subsections.

\begin{figure}[t]
  \centering
  \includesvg[width=\linewidth]{figures/overview.svg}
  \caption{Overview of PAMR. Given a query $(s, r, ?)$, frozen KGE embeddings are retrieved for entities and relations. A bidirectional BFS extracts shortest paths between the subject and each candidate object under the modality-as-entity setting. Each path is encoded as an interleaved sequence of entity and relation tokens, prepended with a CLS token and the query relation. The Transformer encoder contextualises the sequence, and a two-layer MLP readout produces a scalar score per path. Paired positive and negative paths are scored jointly and trained with a margin-based ranking loss.}
  \label{fig:overview}
\end{figure}

\subsection{Embedding Module}
\label{sec:embedding-module}

To supply pre-encoded relational structure without incurring additional
training cost, we use a frozen pre-trained KGE model.  A KGE model
learns entity embeddings $\mathbf{E} \in \mathbb{H}^{|\mathcal{V}|}$
and relation embeddings $\mathbf{R} \in \mathbb{H}^{|\mathcal{R}|}$
such that a scoring function $h: \mathcal{V} \times \mathcal{R} \times
\mathcal{V} \to \mathbb{R}$ assigns higher values to true triples
than to false ones.  The embedding space $\mathbb{H}$ is a
$D$-dimensional continuous manifold.

In a multimodal KG, entities possess heterogeneous modalities.
VL-KGE~\cite{efthymiou2026vlkge} constructs a unified embedding for
every entity by fusing its available modalities into a single point
in $\mathbb{H}$.  For each entity $v \in \mathcal{V}$, let
$\mathcal{M}(v) \subseteq \mathbb{M}$ denote the set of available
modalities.  The unified entity embedding $\mathbf{v}$ is obtained
by averaging over available modalities:
$\mathbf{v} = \frac{1}{|\mathcal{M}(v)|} \sum_{i \in \mathcal{M}(v)} \mathbf{v}_i$,
where $\mathbf{v}_i$ is the modality-specific embedding for modality
$i \in \mathbb{M}$.  This per-entity averaging naturally handles
modality asymmetry, where different entity types carry different
modality combinations without requiring a fixed global fusion scheme.

The effective dimensionality of entity and relation representations
depends on the KGE architecture: complex-valued models output $2D$
dimensions while real-valued models output $D$. The Transformer
hidden dimension is set to $d = \max(d_{\mathcal{V}},
d_{\mathcal{R}})$, and two learnable linear projections align entity
and relation token dimensions to $d$. These
projections are learned during PAMR training while the KGE model
remains frozen.

\subsection{Path Extraction}
\label{sec:path-extraction}

To provide the Transformer with explicit, interpretable connectivity evidence rather than
relying solely on the implicit proximity captured by KGE embeddings, we extract shortest
paths between query pairs via bidirectional breadth-first search (BFS), bounded by a
maximum hop count $K$.

\paragraph{Multimodal and non-multimodal configurations.}
We define two graph configurations that produce qualitatively different path spaces.
The multimodal configuration $\mathcal{G}_{\text{multi}}$ comprises the complete
graph with all entity types, where paths may bridge across different entity types,
exploiting cross-type relationships for reasoning.  The concept-only configuration
$\mathcal{G}_{\text{concept}}$ restricts the graph to conceptual entities and
edges where both endpoints are conceptual, confining paths to a homogeneous
subgraph that excludes multimodal information.  The latter corresponds to
Path-Aware Reasoning (PAR), which relies solely on structural connectivity
without multimodal signals.

Formally, a path $p$ between $s$ and $\hat{o}$ comprises $\ell + 1$ edges, where
$\ell$ denotes the hop count bounded by $K$.
Paths with $\ell \leq 0$ are excluded from training and evaluation as they
carry no structural reasoning signal, which include unreachable pairs and
trivial self loops.

\subsection{PAMR Architecture}
\label{sec:pamr-architecture}

\subsubsection{Token Interleaving}
\label{sec:token-interleaving}

To encode path structure as a flat sequence, entity and relation tokens are interleaved
with a classification token and the query relation prepended.
Let $\mathbf{v}_j \in \mathbb{R}^{d}$ be the entity token for entity $j$ in the path,
obtained from the unified entity embedding via a learnable projection.
A path with $\ell$ hops comprises $\ell+1$ entity tokens and $\ell$ edge tokens.
These are interleaved into a token sequence of length $L = 2\ell + 3$:
\begin{equation}
  \mathbf{T} = \bigl[\mathbf{c};\; \mathbf{r};\;
    \mathbf{v}_{s};\; \mathbf{r}_{1};\; \mathbf{v}_{2};\; \mathbf{r}_{2};\;
  \mathbf{v}_{3};\; \ldots;\; \mathbf{v}_{o}\bigr]
  \in \mathbb{R}^{L \times d}
  \label{eq:token-seq}
\end{equation}
Here $\mathbf{c} \in \mathbb{R}^{d}$ is a learnable classification token
at position 0, initialised from $\mathcal{N}(0, 0.02^2)$.
$\mathbf{r}$ is the query relation representation at position 1, while entity tokens
$\mathbf{v}_{s}, \mathbf{v}_{2}, \ldots, \mathbf{v}_{o}$ occupy positions $2 + 2i$
for $i = 0, \ldots, \ell$, and edge tokens $\mathbf{r}_{1}, \ldots, \mathbf{r}_{\ell}$
occupy positions $3 + 2i$ for $i = 0, \ldots, \ell - 1$.
All positions receive absolute positional embeddings
$\mathbf{p}_i = \operatorname{PE}(i)$ via a learned embedding
$\operatorname{PE}: \{0, \ldots, 2K+2\} \to \mathbb{R}^{d}$,
where $K$ is the maximum hop count.
The final token input to the Transformer is $\mathbf{T}^{\text{in}}_i = \mathbf{T}_i \;+\; \mathbf{p}_i$
\subsubsection{Transformer Encoder}

\label{sec:transformer-encoder}

The token sequence is contextualised through a stack of $L$
Transformer encoder layers~\cite{vaswani2017attention} with $h$ attention heads each:
\begin{equation}
  \mathbf{H}^{(0)} = \mathbf{T}^{\text{in}},\qquad
  \mathbf{H}^{(l)} = \operatorname{EncLayer}^{(l)}\bigl(\mathbf{H}^{(l-1)},\; \mathbf{M}\bigr)
  \quad (l = 1, \ldots, L)
  \label{eq:transformer}
\end{equation}
Here $\mathbf{M} \in \{0, 1\}^{L}$ is a padding mask ($\mathbf{M}_i = 1$ indicates an
invalid position) for paths shorter than the batch maximum $N_{\max}(\mathcal{B})$.
Each encoder layer applies:
\begin{align}
  \tilde{\mathbf{H}} &= \operatorname{LayerNorm}\bigl(
  \mathbf{H}^{(l-1)} + \operatorname{MHA}(\mathbf{H}^{(l-1)}; \mathbf{M})\bigr)
  \label{eq:attn}\\
  \mathbf{H}^{(l)} &= \operatorname{LayerNorm}\bigl(
  \tilde{\mathbf{H}} + \operatorname{FFN}(\tilde{\mathbf{H}})\bigr)
  \label{eq:ffn}
\end{align}
with $\operatorname{FFN}(\mathbf{x}) = \mathbf{W}^{(2)} \operatorname{GELU}(\mathbf{W}^{(1)} \mathbf{x} + \mathbf{b}^{(1)}) + \mathbf{b}^{(2)}$
and inner dimension $d' = 4 \times d$.
GELU \cite{hendrycks2016gaussian} activation
and Post-LN layer normalisation are used throughout.

\subsubsection{Readout and Scoring}

\label{sec:readout}

To obtain a fixed-dimensional path representation, the contextualised token
sequence is aggregated from
$\mathbf{H}^{(L)} \in \mathbb{R}^{L \times d}$
via one of two readout strategies: CLS extraction, which takes the
classification token at position 0, and mean pooling, which averages
over valid tokens.
A two-layer MLP maps the aggregated representation to a scalar score
$\hat{y} = \mathbf{w}_2^{\top} \operatorname{GELU}(\mathbf{W}_1 \mathbf{z} + \mathbf{b}_1) + b_2$
with dropout $p = 0.1$ applied after GELU.
All linear layers use Xavier
uniform initialisation and zero bias.

\subsection{Training}

\label{sec:training}

To train a path scorer with stable supervision, each query must contribute equal gradient
signal regardless of its number of negative samples.
We construct a paired sampling
dataset that guarantees both balance and fairness.
\subsubsection{Paired Sampling}
\label{sec:paired-sampling}

For each query $q$, the query generation process produces one
\textit{positive} path $p^+_q$, which connects $s$ to the
true object $o$, and $N^-$ \textit{negative} paths
$\mathcal{S}^-(q) = \{p^-_{q,1}, \ldots, p^-_{q,N^-}\}$, which connect
$s$ to sampled false objects. The total number of pairs equals
$\sum_q |\mathcal{S}^-(q)|$, ensuring that every negative is visited exactly
once per epoch. Collation interleaves positive and negative pairs
into flat sequences
$[p^+_0, p^-_0, \ldots, p^+_{B-1}, p^-_{B-1}]$.
Here $B$ is the query batch size. Each training batch contains
exactly $B$ positive and $B$ negative samples, producing a guaranteed
balanced composition regardless of the underlying positive to negative
ratio. Variable length paths are zero-padded to the batch maximum
$N_{\max}(\mathcal{B}) = \max_{p \in \mathcal{B}} (\ell_p + 1)$
entity tokens, with padding masks $\mathbf{M}$ as defined in
Section~\ref{sec:token-interleaving}.

\subsubsection{Loss Function}
\label{sec:loss-functions}

PAMR is trained with a margin-based ranking loss that groups samples
by $(s, r, ?)$ and enforces a margin between the positive score
and all negative scores for that query. Let $\mathcal{B} = \{(\hat{y}_1, y_1), \ldots, (\hat{y}_{2B}, y_{2B})\}$
denote a training batch with $y_i \in \{0, 1\}$. The loss is defined as:
\begin{equation}
  \mathcal{L}_{\text{rank}} = \frac{1}{|\mathcal{Q}(\mathcal{B})|} \sum_{q \in \mathcal{Q}(\mathcal{B})}\; \frac{1}{|S^+(q)||S^-(q)|} \sum_{\hat{y}^+ \in S^+(q)} \sum_{\hat{y}^- \in S^-(q)}
  \max\bigl(0,\, m - \hat{y}^+ + \hat{y}^-\bigr)
  \label{eq:loss-rank}
\end{equation}
Here $\mathcal{Q}(\mathcal{B})$ denotes the set of unique queries in $\mathcal{B}$, and
$m = 1.0$ is the margin. This loss directly optimises the ranking objective and prevents
score collapse.

\section{Experiment}
\label{sec:experiment}
\subsection{Implementation Details}
\label{sec:implementation}

All experiments run on a Vast.ai instance equipped with an NVIDIA
GeForce RTX 5060 Ti and an AMD EPYC 7763 processor allocated with 8 threads
and 16\,GB virtual RAM. The software stack comprises Python~3.11,
PyTorch~2.7.1 with CUDA~12.8 and cuDNN~9.

\paragraph{Path extraction.}
The bidirectional BFS of Section~\ref{sec:path-extraction} is
implemented in C++ (GNU GCC 13.2, compiled with C++23 standards)
and invoked from Python. The graph adjacency list is loaded once
into memory, and queries are distributed across 8 threads for
parallel path search. Per-thread outputs are merged into a single
path table per split. Precomputed paths are cached on disk as
pickle files for fast reuse across training runs.

\paragraph{GPU memory management.}
To accommodate the large frozen embedding table alongside the trainable
Transformer, GPU memory is managed through expandable memory segments
and periodic cache defragmentation after model initialisation.

\paragraph{Paired sampling implementation.}
The paired dataset is constructed once per epoch from precomputed
path tables, with per-query pairing parallelised across CPU cores.
Variable length paths are zero-padded to the batch maximum and
collated with padding masks as described in
Section~\ref{sec:paired-sampling}.

\paragraph{Training configuration.}
PAMR is trained with AdamW~\cite{loshchilov2019decoupled} at learning
rate $\eta = 0.01$, weight decay $\lambda = 0.01$, $\beta_1 = 0.9$,
$\beta_2 = 0.999$. The ranking loss (Eq.~\ref{eq:loss-rank}) is
used with margin $m = 1.0$. Batch size is 8 queries (16 samples per
batch with the guaranteed balanced positive/negative composition).
Training proceeds for 50 epochs. The Transformer encoder uses
$L = 8$ layers and $h = 8$ attention heads with
$d' = 2048$ and dropout $p = 0.1$. The hidden dimension
$d = 1536$ is auto-detected from the frozen VL-KGE
model (ComplEx, $D = 768$). Only PAMR parameters are updated; the
frozen VL-KGE model remains untouched.

\subsection{Datasets}
\label{sec:datasets}

Experiments are conducted on WikiArt-MKG~v2~\cite{efthymiou2026vlkge},
a large-scale multimodal knowledge graph of fine art containing
224,166 entities and 22 relation types. Artworks are represented
visually through image embeddings, while attributes such as styles,
genres, years, and locations are represented textually. This modality
asymmetry requires models to handle entities with different
combinations of available modalities. The dataset comprises 224,166
entities, 216,564 artworks and 7,602 attributes, connected by 22
relation types, totalling 7,877,220 training triples, 208,513
validation triples, and 208,368 test triples. The graph is split
chronologically: training data precedes validation data, which in
turn precedes test data. Negative samples are generated via rejection
sampling with 1 negative per training query and 10 negatives per
validation/test query. For evaluation, three relation types
introduced as preprocessing inverses
and 10 negatives per validation/test query. For evaluation, three
relation types introduced as preprocessing inverses
(\texttt{isRelatedToArtwork}, \texttt{isRelatedToArtist},
\texttt{hasCreatedArtwork}) are excluded, leaving the four
relations between artists (\texttt{isInfluencedBy},
\texttt{isInfluencedOn}, \texttt{isPupilOf}, \texttt{isTeacherOf})
as primary evaluation targets. These four relations are chosen
because they connect artist entities, enabling evaluation under both
multimodal and non-multimodal reasoning settings.

Path extraction uses bidirectional BFS with max hops $H = 5$, base
fanout $F = 100$, and decay factor $\lambda = 0.5$
\cite{icaart26}. Both multimodal and non-multimodal configurations are evaluated, providing a controlled comparison that isolates the contribution of multimodal path structure.

\subsection{Evaluation Metrics}
\label{sec:evaluation-metrics}

For each query, the model assigns a scalar score to every candidate
object and ranks them in descending order. We report Mean Reciprocal
Rank (MRR) and Hit-rate at cutoff $K$ (Hits@$K$), computed both
globally and disaggregated by relation type.

\subsection{Results}
\label{sec:results}

For all the conducted experiments, baselines built upon VL-KGE~\cite{efthymiou2026vlkge} are used with different embedding methods, including TransE, RotatE, DistMult, and ComplEx. We explicitly did not experiment with any LLM-based methods due to them being pre-trained on vast amount of external corpora which makes comparisons unfair.

\subsubsection{Comparative Study}

Table~\ref{tab:results} reports macro-averaged MRR and Hits@$K$ across
the four evaluation relations between artists. All methods share the same
frozen ComplEx embeddings to
ensure a fair comparison. The VL-KGE baselines score triples directly
without path reasoning, while PAR uses only structural paths and PAMR
additionally exploits multimodal token representations.

\begin{table*}[t]
  \centering
  \setlength{\tabcolsep}{2pt}
  \renewcommand{\arraystretch}{1.15}
  {\footnotesize
    \begin{tabular}{l|cccc|cc}
      \toprule
      \multirow{2}{*}{Metric} &
      \multicolumn{4}{|c|}{VL-KGE Embedding Baselines} &
      \multicolumn{2}{|c|}{Ours} \\
      \cmidrule(lr){2-5} \cmidrule(lr){6-7}
      & TransE & RotatE & DistMult & ComplEx & PAR & PAMR \\
      \midrule
      MRR  & 0.3196 & 0.3272 & 0.3107 & \underline{0.3630} & 0.2212 & \textbf{0.3789} \\
      Hits@1 & 0.2832 & 0.1824 & 0.3181 & \textbf{0.3870} & 0.0576 & \underline{0.3455} \\
      Hits@3 & 0.4938 & 0.4188 & 0.5283 & \textbf{0.6245} & 0.1987 & \underline{0.5759} \\
      Hits@5 & 0.7032 & 0.5973 & 0.7046 & \underline{0.7719} & 0.3306 & \textbf{0.7819} \\
      \bottomrule
  \end{tabular}}
  \caption{Macro-averaged ranking metrics across four evaluation
    relations between artists. \textbf{Bold} indicates the best per metric;
  \underline{underlines} indicate the second best.}
  \label{tab:results}
\end{table*}

The table reveals a consistent ordering across method
families. The results naturally stratify into three tiers: the four
VL-KGE embedding baselines occupy a moderate performance band, PAR
falls below this band, and PAMR surpasses it. Among the baselines,
ComplEx consistently outperforms TransE, RotatE, and DistMult,
achieving 0.363~MRR compared to their 0.31--0.33~MRR range, while the
latter three cluster tightly together, differing by less than
0.02~MRR.

Adding multimodal token representations through PAMR surpasses the
strongest baseline. PAMR achieves the best MRR at 0.379, a 4.4\%
improvement over ComplEx, and the best Hits@5 at 0.782, a 1.3\%
improvement. At the stricter cutoffs, PAMR remains competitive with
ComplEx: Hits@1 reaches 0.345 versus ComplEx's 0.387, and Hits@3
reaches 0.576 versus ComplEx's 0.625. The crossover between MRR and
Hits@5, where PAMR leads, and Hits@1 and Hits@3, where ComplEx leads,
reflects a difference in scoring behaviour: ComplEx concentrates
probability on a small set of candidates, while PAMR distributes
probability more broadly, improving average rank and deeper recall
metrics. This pattern aligns with the design goal of path-aware
reasoning, where explicit structural evidence supports robust ranking
decisions across the entire candidate set. Moreover, PAMR trains in
approximately 543 seconds per epoch compared to 577 seconds for
ComplEx, demonstrating that the path-based Transformer architecture is
computationally competitive with direct triple scoring.

\subsubsection{Per-Relation Analysis}

To complement the macro-averaged results, Figures~\ref{fig:mrr},
\ref{fig:hits1}, \ref{fig:hits3}, and~\ref{fig:hits5} present the
per-relation breakdown for each metric. The four relations exhibit
distinct performance profiles, indicating that the benefit of
multimodal path reasoning depends on the relational context.

\begin{figure*}[t]
  \centering
  \includesvg[width=\linewidth]{figures/MRR.svg}
  \caption{Per-relation MRR across all methods.}
  \label{fig:mrr}
\end{figure*}

\begin{figure*}[t]
  \centering
  \includesvg[width=\linewidth]{figures/Hits1.svg}
  \caption{Per-relation Hits@1 across all methods.}
  \label{fig:hits1}
\end{figure*}

\begin{figure*}[t]
  \centering
  \includesvg[width=\linewidth]{figures/Hits3.svg}
  \caption{Per-relation Hits@3 across all methods.}
  \label{fig:hits3}
\end{figure*}

\begin{figure*}[t]
  \centering
  \includesvg[width=\linewidth]{figures/Hits5.svg}
  \caption{Per-relation Hits@5 across all methods.}
  \label{fig:hits5}
\end{figure*}

Two relations strongly favour PAMR. On \texttt{isPupilOf}, PAMR
achieves 0.390~MRR and 0.909~Hits@5, the highest Hits@5 across all
methods and relations, surpassing ComplEx by 0.038 MRR and
0.091~Hits@5. On \texttt{isInfluencedOn}, PAMR attains 0.415 MRR,
exceeding ComplEx by 0.090 MRR, the largest single-relation gap in
the experiment. These relations connect artists through stylistic
lineage and creative influence, where shared visual motifs and textual
descriptions of artistic movements provide rich multimodal context
that the Transformer exploits when scoring paths.

In contrast, \texttt{isTeacherOf} and \texttt{isInfluencedBy} favour
ComplEx (0.422 and 0.354 MRR versus PAMR's 0.390 and 0.321),
reflecting more abstract associations where direct embedding proximity
suffices. PAR remains the weakest method across all four relations,
confirming that path structure alone, without multimodal tokens,
provides insufficient signal in the concept-only subgraph
$\mathcal{G}_{\text{concept}}$.

\subsection{Case Study}
\label{sec:case-study}

To qualitatively demonstrate the advantages of incorporating multimodal entities into the reasoning process, we present a case study using the query (Pablo Picasso, isInfluenceOn, ?) extracted from the test set as illustrated by Figure \ref{fig:case-study}. We evaluate two candidate answers, Salvador Dali and Lee Krasner, based on the provided subgraph.

\begin{figure}[t]
  \centering
  \includesvg[width=\linewidth]{figures/case.svg}
  \caption{Case Study. A query is given asking for whom being influenced by Picasso between Salvador Dali and Lee Krasner, contextualized with a subgraph containing intermediate artist entities and their structural relations. Besides that, the artworks are also utilized for multimodal information, specifically the Man with a Porron by Salvador Dali, Child with a Dove by Pablo Picasso, and Self-Portrait by Henri Matisse.}
  \label{fig:case-study}
\end{figure}

Under the concept-only configuration used by PAR, the bidirectional BFS extracts paths containing only artist entities. The path to Salvador Dali traverses through intermediate contemporaries, such as Pablo Picasso $\xrightarrow{\text{isRelatedArtist}}$ Juan Gris $\xrightarrow{\text{isRelatedArtist}}$ Salvador Dali. Similarly, the path to Lee Krasner relies on structural artist connections, such as Pablo Picasso $\xrightarrow{\text{isRelatedArtist}}$ Georges Braque $\xrightarrow{\text{isRelatedArtist}}$ Henri Matisse $\xleftarrow{\text{isInfluenceBy}}$ Lee Krasner. These paths merely capture social or professional proximity without providing concrete artistic evidence.

Conversely, under the multimodal-as-entity setting used by PAMR, artworks are treated as individual nodes bridging the artists. The shortest path to Salvador Dali explicitly routes through their creative works, such as Pablo Picasso $\xleftarrow{\text{hasArtist}}$ \textit{Child with a Dove} $\xrightarrow{\text{isRelatedArtwork}}$ \textit{Man with a Porron} $\xrightarrow{\text{hasArtist}}$ Salvador Dali. The path to Lee Krasner similarly routes through an artwork related to Henri Matisse: Pablo Picasso $\xleftarrow{\text{hasArtist}}$ \textit{Child with a Dove} $\xrightarrow{\text{isRelatedArtwork}}$ \textit{Self-Portrait} $\xrightarrow{\text{hasArtist}}$ Henri Matisse $\xleftarrow{\text{isInfluenceBy}}$ Lee Krasner.

When processing these token sequences, the Transformer encoder in PAMR correctly ranks Salvador Dali above Lee Krasner. By projecting the frozen multimodal embeddings of the artworks into the token space, PAMR assesses the actual visual and semantic relationship between the pieces. The explicit presence of connected artworks provides a tangible explanation for the influence, predicting that Picasso had a more direct, verifiable impact on Dali through their related works. Artistically, the path bridging Picasso to Dali via \textit{Man with a Porron} captures a highly cohesive semantic lineage: both artists share deep roots in the Spanish avant-garde, and their multimodal embeddings reflect shared structural experimentation and foundational Cubist motifs. In contrast, the path to Lee Krasner via Matisse's \textit{Self-Portrait} traverses disparate artistic movements, bridging European Cubism to Fauvism, and ultimately to American Abstract Expressionism. The substantial visual and semantic leap required across these distinct, geographically and chronologically distant stylistic paradigms makes the structural path to Krasner far less probable. This highlights how PAMR not only improves ranking accuracy but also yields highly interpretable reasoning trajectories grounded in rich multimodal data and art-historical semantics compared to purely conceptual graphs.

\begin{credits}
  \subsubsection{\ackname} This research is funded by the University of Science, VNU-HCM, Vietnam under grant number CNTT 2025-16.

  \subsubsection{\discintname}
  The authors have no competing interests to declare that are
  relevant to the content of this article.
\end{credits}
%
% ---- Bibliography ----
%
% BibTeX users should specify bibliography style 'splncs04'.
% References will then be sorted and formatted in the correct style.
%
\bibliographystyle{splncs04}
\bibliography{refs}

% \begin{thebibliography}{8}
% \bibitem{ref_article1}
% Author, F.: Article title. Journal \textbf{2}(5), 99--110 (2016)

% \bibitem{ref_lncs1}
% Author, F., Author, S.: Title of a proceedings paper. In: Editor,
% F., Editor, S. (eds.) CONFERENCE 2016, LNCS, vol. 9999, pp. 1--13.
% Springer, Heidelberg (2016). \doi{10.10007/1234567890}

% \bibitem{ref_book1}
% Author, F., Author, S., Author, T.: Book title. 2nd edn. Publisher,
% Location (1999)

% \bibitem{ref_proc1}
% Author, A.-B.: Contribution title. In: 9th International Proceedings
% on Proceedings, pp. 1--2. Publisher, Location (2010)

% \bibitem{ref_url1}
% LNCS Homepage, \url{http://www.springer.com/lncs}, last accessed 2023/10/25
% \end{thebibliography}
\end{document}
